\section{Introduction}
\label{sec:intro}

What does quantum mechanics do to a system whose classical motion is chaotic? Chaos in classical mechanics
means exponential sensitivity to initial conditions: two nearby trajectories separate like $\ee^{\lambda t}$,
where $\lambda$ is the Lyapunov exponent. Quantum mechanics has no trajectories, and the Schr\"odinger equation
is linear, so nothing can separate exponentially in that sense. Yet the classical limit must come out right.
The kicked rotor is the simplest system in which this tension can be studied in full detail, and it is the
subject of this project.

The kicked rotor is a particle on a ring (a rigid rotor) that receives a sharp kick once per period, with a
strength that depends on its angle $\theta$:
\begin{equation}
  H = \frac{p^2}{2} + K\cos\theta\,\sum_{n=-\infty}^{\infty}\delta(t - n).
  \label{eq:hamiltonian}
\end{equation}
Between kicks the rotor turns freely. Integrating Hamilton's equations over one period gives the Chirikov
standard map \cite{Chirikov1979}, one of the most studied models of Hamiltonian chaos. For weak kicks most
orbits are regular and the momentum stays bounded. Above a critical kick strength
$K_c \approx 0.9716$ \cite{Greene1979} the last barrier across the momentum direction breaks, and for
$K \gtrsim 5$ almost all orbits are chaotic. The momentum then performs a random walk: averaged over an
ensemble, $\avg{p^2}$ grows linearly in time, $\avg{p^2} \approx Dn$ after $n$ kicks, with a diffusion rate
$D \approx K^2/2$ plus corrections that oscillate with $K$ \cite{Chirikov1979,Rechester1980,Rechester1981}.

The quantum rotor behaves very differently. Casati, Chirikov, Izraelev and Ford \cite{Casati1979} solved it
numerically in 1979. At first the quantum energy grew like the classical one, but after a break time the
growth slowed down. Later work showed that it stops altogether: the momentum distribution freezes into an
exponentially localized profile $|\psi_m|^2 \sim \ee^{-2|m|/\xi}$ around its starting point. This is
\emph{dynamical localization}. Fishman, Grempel and Prange \cite{Fishman1982} explained it by mapping the
eigenvalue problem of the kicked rotor onto a one-dimensional lattice with pseudo-random on-site energies,
the Anderson model, in which every eigenstate is exponentially localized. The localization length is set by
the classical diffusion rate, $\xi \approx D/\hbar^2$ in our units \cite{Shepelyansky1986,Izrailev1990}. The
effect was observed directly in 1995 with ultracold sodium atoms in a pulsed standing wave of light
\cite{Moore1995}.

In this project we simulate both the classical and the quantum rotor, check the code against exact results,
and then use it to study chaos, diffusion and dynamical localization. The main result is
\cref{fig:quantum_vs_classical}: the quantum energy follows the classical diffusive growth for a few kicks
and then saturates, at a level and after a time that both scale with $\xi$.

The report is organized as follows. \Cref{sec:model} sets up the model, the quantum one-period
(Floquet) operator and the numerical method, and summarizes the validation. \Cref{sec:classical} covers the
classical dynamics: the transition to chaos and the diffusion rate with its Bessel-function corrections.
\Cref{sec:quantum} is the core: Floquet theory, quantum resonance, the break time, the localized momentum
distribution and its localization length, how $\xi$ scales with $K$ and $\hbar$, and the Anderson
analogy. \Cref{sec:correspondence} discusses quantum-classical correspondence and time reversal, and
\cref{sec:experiments} the cold-atom experiment. The appendices contain the derivations of the classical
correlations and of the exact two-kick results, details of the validation tests, and instructions for
reproducing every figure.
